\subsection{Pfaffian point process}
In this paper, we assume that $\mfrak{X}$ is a countable set. The collection of point configurations in~$\mfrak{X}$ is identified with $\Omega=\Omega (\mfrak{X})=\{0,1\}^{\mfrak{X}}$, which is equipped with the product topology to be a~compact topological space. We regard each element $\omega\in \Omega$ as a function $\omega\colon \mfrak{X}\to\{0,1\}$ or a~collection of points $\omega=\{x_{i}\in\mfrak{X}\}_{i}$. We adopt the $\sigma$-algebra of Borel sets $\Sigma$. Then, for distinct points $x_{1},\dots, x_{n}\in \mfrak{X}$, the cylinder set
\begin{equation*}
	\Omega_{x_{1},\dots, x_{n}}=\{\omega\in\Omega\,|\,\omega (x_{1})=\cdots =\omega (x_{n})=1\}
\end{equation*}
is measurable.
Given a probability measure $M$ on $(\Omega,\Sigma)$, the $n$-point correlation function $\rho^{M}_{n}$, $n\in\mbb{N}$ is an $n$-variable symmetric function defined as the probability weight of cylinder sets:
\begin{equation*}
	\rho^{M}_{n}(x_{1},\dots, x_{n})=M(\Omega_{x_{1},\dots, x_{n}}),
\end{equation*}
where $x_{1},\dots, x_{n}\in \mfrak{X}$ are distinct. It is conventional to extend $\rho^{M}_{n}$ to a function on $\mfrak{X}^{n}$ so that it vanishes if any two points coincide.
Note that a system of correlation functions $\big\{\rho^{M}_{n}\big\}_{n\in\mbb{N}}$ determines the probability measure $M$ uniquely since the cylinder sets generate $\Sigma$.
A random variable $X$ with values in $(\Omega,\Sigma)$ is called a point process in $\mfrak{X}$. We also call a probability measure on $(\Omega,\Sigma)$ a point process not distinguishing a random variable from its distribution.

To define a Pfaffian point process, we need to fix some notations.
Suppose that a $(2\times 2)$-matrix-valued function $\mbb{K}(\cdot,\cdot)\colon \mfrak{X}\times\mfrak{X}\to M(2;\mbb{C})$ satisfying the anti-symmetry
\begin{equation}\label{eq:anti-sym_matrix_func}
	\mbb{K}(x,y)^{\mrm{T}}=-\mbb{K}(y,x),\qquad x, y\in\mfrak{X}
\end{equation}
is given. For any $n\in\mbb{N}$, we adopt the following identification of vector spaces:
\begin{align*}
	M(2;\mbb{C})\otimes M(n;\mbb{C}) & \xrightarrow{\sim} M(2n;\mbb{C}),\\
	 e_{i,j}\otimes e_{k,l}&\mapsto e_{2(k-1)+i, 2(l-1)+j},\qquad i,j=1,2,\quad k,l=1,\dots, n,
\end{align*}
where $e_{i,j}$ is the matrix with entry $1$ at the $(i,j)$-position and $0$ elsewhere.
Given points $x_{1},\dots, x_{n}\in\mfrak{X}$, we understand $[\mbb{K}(x_{i},x_{j})]_{1\le i,j\le n}$ as a $(2n\times 2n)$-matrix under the above identification and by regarding $x_{1},\dots, x_{n}$ as legs for $M(n;\mbb{C})$. The anti-symmetry (\ref{eq:anti-sym_matrix_func}) of the function $\mbb{K}(\cdot,\cdot)$ ensures the anti-symmetry of the $(2n\times 2n)$-matrix $[\mbb{K}(x_{i},x_{j})]_{1\le i,j\le n}$ so that
\begin{equation*}
[\mbb{K}(x_{i},x_{j})]_{1\le i,j\le n}^{\mrm{T}}=\big[\mbb{K}(x_{j},x_{i})^{\mrm{T}}\big]_{1\le i,j\le n}=-[\mbb{K}(x_{i},x_{j})]_{1\le i,j\le n}.
\end{equation*}
Therefore, the Pfaffian $\pf [\mbb{K}(x_{i},x_{j})]_{1\le i,j\le n}$ is defined.

\begin{Definition}
A probability measure $M$ on $(\Omega,\Sigma)$ is a Pfaffian point process (PfPP) if there exists a $(2\times 2)$-matrix-valued function $\mbb{K}^{M}(\cdot,\cdot)\colon \mfrak{X}\times\mfrak{X}\to M(2;\mbb{C})$ satisfying the anti-symmetry~(\ref{eq:anti-sym_matrix_func}) such that every $n$-point function with $n\in\mbb{N}$ admits a Pfaffian expression
\begin{equation*}
	\rho_{n}^{M}(x_{1},\dots, x_{n})=\pf \big[\mbb{K}^{M}(x_{i},x_{j})\big]_{1\le i,j\le n}.
\end{equation*}
We call the matrix-valued function $\mbb{K}^{M}(\cdot,\cdot)$ a correlation kernel of the PfPP $M$.
\end{Definition}

For a PfPP $M$, we write its correlation kernel as
\begin{equation*}
	\mbb{K}^{M}(x,y)
	=\left(
	\begin{matrix}
		\mbb{K}_{11}^{M}(x,y) & \mbb{K}_{12}^{M}(x,y) \\
		\mbb{K}_{21}^{M}(x,y) & \mbb{K}_{22}^{M}(x,y)
	\end{matrix}
	\right),\qquad x,y\in\mfrak{X}.
\end{equation*}
If, in particular, $\mbb{K}_{11}^{M}(x,y)=\mbb{K}_{22}^{M}(x,y)=0$ identically holds, we have
\begin{equation*}
	\mrm{Pf}\big[\mbb{K}^{M}(x_{i},x_{j})\big]_{1\le i,j\le n}=\det \big[\mbb{K}_{12}^{M}(x_{i},x_{j})\big]_{1\le i,j\le n}.
\end{equation*}
Therefore, we may naturally have the following definition as a special case of PfPPs.

\begin{Definition}
A PfPP $M$ on $(\Omega,\Sigma)$ is called a determinantal point process (DPP) if $\mbb{K}^{M}_{11}(x,y)\allowbreak =\mbb{K}^{M}_{22}(x,y)=0$ identically.
In this case, each correlation function admits a determinantal expression
\begin{equation*}
	\rho_{n}^{M}(x_{1},\dots, x_{n})=\det \big[K^{M}(x_{i},x_{j})\big]_{1\le i,j\le n},\qquad K^{M}(x,y)=\mbb{K}^{M}_{12}(x,y),\qquad x,y\in\mfrak{X}
\end{equation*}
We adopt the abuse of terminology to call the function $K^{M}(\cdot,\cdot)$ on $\mfrak{X}\times\mfrak{X}$ a correlation kernel of the DPP $M$.
\end{Definition}

DPPs form a significant class of point processes that have many applications to random partition and asymptotic representation theory \cite{Borodin1998a, BorodinOlshanski1998a, Borodin1998b,BorodinOlshanski1998b,Okounkov2001,Olshanski1998, Olshanski2003} (see also \cite{Borodin2011,BenHoughKrishnapurPeresVirag2006, Lyons2003, Soshnikov2000}). The analogy between DPPs and fermionic calculi is well-known \cite{Lytvynov2002,LytvynovMei2007,Olshanski2020, ShiraiTakahashi2003a, ShiraiTakahashi2003b, Spohn1987}. In particular, a quasi-free state of a certain class over an algebra of anti-commutation relations gives a DPP.
Compared to DPPs, there seem to be less attempts to unify an interplay between PfPPs and free fermions, while there are many preceding studies on concrete examples (incomplete references include \cite{BorodinRains2005, Ferrari2004, KatoriTanemura2007,Matsumoto2005, MatsumotoShirai2013,Nagao2007,Rains2000,Vuletic2007,WangLi2019}).
One of the aims of this paper is to extend a portion of the insight on DPPs found in the above mentioned literatures and to establish a framework to study PfPPs from free fermionic perspectives. Many settings, notations and problems are motivated by \cite{Olshanski2020}.

\subsection{From a positive contraction to a PfPP}
A well-known construction of a DPP on $\mfrak{X}$ \cite{Soshnikov2000} starts from an operator $K$ on $\ell^{2}(\mfrak{X})$ such that $K=K^{\ast}$ and $0\le K\le 1$, namely, a~positive contraction on~$\ell^{2}(\mfrak{X})$. Given a~positive contraction~$K$ on $\ell^{2}(\mfrak{X})$, there exists a DPP $M^{K}$ on $\mfrak{X}$ such that each correlation function is given by
\begin{equation*}
	\rho^{M^{K}}_{n}(x_{1},\dots, x_{n})=\det [K(x_{i},x_{j}) ]_{1\le i,j\le n},\quad K(x,y)=(e_{x},Ke_{y})_{\ell^{2}(\mfrak{X})}.
\end{equation*}
Here $e_{x}\in \ell^{2}(\mfrak{X})$, $x\in\mfrak{X}$ is a function defined by $e_{x}(y)=\delta_{x,y}$, $y\in\mfrak{X}$. Then the collection $\{e_{x}\}_{x\in\mfrak{X}}$ forms a complete orthonormal system of $\ell^{2}(\mfrak{X})$.

Here, we give a direct generalization of this result to PfPPs.
We define the complex conjugate~$J$ of $\ell^{2}(\mfrak{X})$ as an anti-linear operator on it that fixes each function~$e_{x}$, $x\in\mfrak{X}$.
We set $\mcal{K}=\ell^{2}(\mfrak{X})\oplus \ell^{2}(\mfrak{X})$ and take an anti-unitary involution $\Gamma$ on $\mcal{K}$ defined by
\begin{equation}
\label{eq:Gamma}
	\Gamma=\left(
	\begin{matrix}
	0 & J \\
	J & 0
	\end{matrix}
	\right)
\end{equation}
according to the prescribed direct sum decomposition.
For this pair $(\mcal{K},\Gamma)$, we consider the collection of operators
\begin{equation}
\label{eq:collection_covariance_ops}
	\mcal{Q}(\mcal{K},\Gamma)=\big\{S\in\mfrak{B}(\mcal{K})\,|\,0\le S=S^{\ast}\le 1,\, \overline{S}=1-S\big\},
\end{equation}
where $\mfrak{B}(\mcal{K})$ is the set of bounded operators on $\mcal{K}$, and for any operator $A\in\mfrak{B}(\mcal{K})$, we write $\overline{A}:=\Gamma A\Gamma$.
We will show that each operator $S\in\mcal{Q}(\mcal{K},\Gamma)$ uniquely determines a PfPP.

\begin{Proposition}\label{prop:existence_ppp}
Let us take $S\in\mcal{Q}(\mcal{K},\Gamma)$ and write it as
\begin{equation*}
	S=\left(
	\begin{matrix}
		S_{11} & S_{12} \\
		S_{21} & S_{22}
	\end{matrix}
	\right).
\end{equation*}
There exists a PfPP $M^{S}$ on $\mfrak{X}$ such that each correlation function is given by
\begin{equation*}
	\rho^{M^{S}}_{n}(x_{1},\dots, x_{n})=\mrm{Pf} [\mbb{K}_{S}(x_{i},x_{j}) ]_{1\le i,j\le n},
\end{equation*}
where
\begin{equation}
\label{eq:matrix_kernel_S}
	\mbb{K}_{S}(x,y)=\left(
	\begin{matrix}
		(e_{x},S_{21}e_{y})_{\ell^{2}(\mfrak{X})} & (e_{x},S_{22}e_{y})_{\ell^{2}(\mfrak{X})} \\
		(e_{x},(S_{11}-1)e_{y})_{\ell^{2}(\mfrak{X})} & (e_{x},S_{12}e_{y})_{\ell^{2}(\mfrak{X})}
	\end{matrix}
	\right),\qquad x,y\in\mfrak{X}.
\end{equation}
\end{Proposition}

\begin{Remark}
If $S$ is diagonal; $S_{12}=S_{21}=0$, the PfPP $M^{S}$ is a DPP.
\end{Remark}

\begin{Remark}
\label{rem:check_anti-sym_matrix_func}
For $S\in\mcal{Q}(\mcal{K},\Gamma)$, the function $\mbb{K}_{S}(\cdot,\cdot)$ satisfies the required anti-symmetry (\ref{eq:anti-sym_matrix_func}). In fact, the self-adjointness of $S$ gives $S_{11}^{\ast}=S_{11}$, $S_{22}^{\ast}=S_{22}$, $S_{12}^{\ast}=S_{21}$, and the identity $\overline{S}=1-S$ implies $JS_{11}J=1-S_{22}$, $JS_{12}J=-S_{21}$. In particular, we have $S_{12}^{\mrm{T}}=-S_{12}$, $S_{21}^{\mrm{T}}=-S_{21}$. Therefore,
\begin{gather*}
	(e_{x},S_{21}e_{y})_{\ell^{2}(\mfrak{X})} =-(e_{y},S_{21}e_{x})_{\ell^{2}(\mfrak{X})},\\
	(e_{x},S_{12}e_{y})_{\ell^{2}(\mfrak{X})} =-(e_{y},S_{12}e_{x})_{\ell^{2}(\mfrak{X})}, \\
	(e_{x},S_{22}e_{y})_{\ell^{2}(\mfrak{X})} =\overline{(e_{y},J(1-S_{11})Je_{x})_{\ell^{2}(\mfrak{X})}}=-(e_{y},(S_{11}-1)e_{x})_{\ell^{2}(\mfrak{X})}, \\
(e_{x},(S_{11}-1)e_{y})_{\ell^{2}(\mfrak{X})}=-\overline{(e_{y},JS_{22}Je_{x})_{\ell^{2}(\mfrak{X})}}=-(e_{y},S_{22}e_{x})_{\ell^{2}(\mfrak{X})}
\end{gather*}
for any $x,y\in\mfrak{X}$, which implies the anti-symmetry (\ref{eq:anti-sym_matrix_func}).
\end{Remark}
